\documentclass[12pt]{article}
\usepackage{amsmath,amssymb,amsthm}
\usepackage[margin=1in]{geometry}
\usepackage{algorithm,algpseudocode}

\newtheorem{theorem}{Theorem}
\newtheorem{lemma}[theorem]{Lemma}
\newtheorem{corollary}[theorem]{Corollary}
\newtheorem{proposition}[theorem]{Proposition}
\theoremstyle{definition}
\newtheorem{definition}[theorem]{Definition}

\title{Project Euler Problem 429:\\Sum of Squares of Unitary Divisors}
\author{}
\date{}

\begin{document}
\maketitle

\section{Problem Statement}
A unitary divisor $d$ of a number $n$ is a divisor such that $\gcd(d, n/d) = 1$.
Find $S(100\,000\,000!) \pmod{1\,000\,000\,009}$, where $S(n)$ is the sum of the squares of the unitary divisors of $n$.

\section{Mathematical Derivation}

\subsection{Unitary Divisor Sum Formula}
\begin{theorem}
If $n = p_1^{a_1} p_2^{a_2} \cdots p_k^{a_k}$, then the sum of squares of unitary divisors of $n$ is:
\[
\sigma_2^*(n) = \prod_{i=1}^{k} \left(1 + p_i^{2a_i}\right)
\]
\end{theorem}

\begin{proof}
A unitary divisor $d$ of $n$ satisfies $\gcd(d, n/d) = 1$. For each prime power $p_i^{a_i}$ dividing $n$, the divisor $d$ either contains $p_i^{a_i}$ or $p_i^0 = 1$ (no intermediate powers, as those would share a common factor with $n/d$). Thus:
\[
\sigma_2^*(n) = \sum_{\substack{d \mid n \\ \gcd(d, n/d) = 1}} d^2 = \prod_{i=1}^{k} \left(1 + (p_i^{a_i})^2\right) = \prod_{i=1}^{k} \left(1 + p_i^{2a_i}\right)
\]
\end{proof}

\subsection{Legendre's Formula}
\begin{theorem}[Legendre]
The exponent of prime $p$ in $N!$ is:
\[
v_p(N!) = \sum_{i=1}^{\infty} \left\lfloor \frac{N}{p^i} \right\rfloor
\]
\end{theorem}

\subsection{Solution}
For $N = 10^8$, we compute:
\[
S(N!) = \prod_{p \leq N} \left(1 + p^{2 \cdot v_p(N!)}\right) \pmod{10^9 + 9}
\]
using modular exponentiation. Since $10^9 + 9$ is prime, we can use Fermat's little theorem if needed.

\section{Result}
\[
\boxed{S(100\,000\,000!) \equiv 98792821 \pmod{1\,000\,000\,009}}
\]

\section{Answer}

\[
\boxed{98792821}
\]

\end{document}
